\documentclass[aps,prd,reprint,nofootinbib,notitlepage]{revtex4-2}
\usepackage[T1]{fontenc}
\usepackage{amsmath,amssymb,bm,booktabs}
\usepackage[colorlinks=true,allcolors=blue]{hyperref}

\newcommand{\nue}{\nu_e}
\newcommand{\anue}{\bar{\nu}_e}
\newcommand{\nux}{\nu_x}
\newcommand{\dd}{\mathrm{d}}
\newcommand{\eq}{\mathrm{eq}}

\begin{document}

\title{WeakLib neutrino interactions in the grey M1 formulation of nuX}
\date{}

\begin{abstract}
We describe the coupling of WeakLib neutrino interaction tables to the grey
M1 radiation equations in nuX. The matter equation of state is supplied by
EOSX, whereas WeakLib supplies energy-dependent absorption and scattering
data. The coupling includes charged-current emission and absorption,
elastic scattering, electron--positron pair annihilation, nucleon--nucleon
bremsstrahlung, and inelastic scattering on electrons and positrons.
Plasmon decay can be added through a separate analytic approximation.
A Fermi--Dirac spectrum inferred from the evolved number and energy
densities is used to evaluate the collision integrals. We distinguish this
spectral approximation from the Minerbo angular closure, explain the
conditions for detailed balance and lepton-number conservation, and compare
the reaction content with WeakRates and NuRates. A final section describes
table preparation, interpolation, and the parameters needed to use the
coupling.
\end{abstract}

\maketitle

\section{Radiation and matter variables}

Grey M1 transport evolves the energy density and flux integrated over
neutrino energy, together with a neutrino-number density. It does not evolve
an independent radiation field at each neutrino energy
\cite{MusolinoRezzolla2023}. The three species groups considered here are
$\nue$, $\anue$, and $\nux$, where the last group is the sum of
$\nu_\mu$, $\bar{\nu}_\mu$, $\nu_\tau$, and $\bar{\nu}_\tau$.
Muon charged-current reactions and neutrino flavor conversion are not
included.

With metric signature $(-,+,+,+)$ and fluid four-velocity $u^\mu$, the
radiation stress-energy tensor is decomposed as
\begin{equation}
 T_s^{\mu\nu}=J_s u^\mu u^\nu
       +H_s^\mu u^\nu+u^\mu H_s^\nu+K_s^{\mu\nu}.
 \label{eq:stress}
\end{equation}
Here $J_s$ is the comoving energy density, $H_s^\mu u_\mu=0$, and
$K_s^{\mu\nu}u_\nu=0$. The comoving number density is denoted by $n_s$.
These are physical local densities, not coordinate-volume-weighted
evolution variables. In units with $c=1$, the exchange equations are
\begin{align}
 \nabla_\mu T_s^{\mu\nu}&=G_s^\nu, &
 \nabla_\mu {\cal N}_s^\mu&=R_s, \label{eq:radiation_exchange}\\
 \nabla_\mu T_{\mathrm{m}}^{\mu\nu}&=-\sum_s G_s^\nu, &
 \nabla_\mu(n_bY_eu^\mu)&=-R_{\nue}+R_{\anue}.
 \label{eq:matter_exchange}
\end{align}
The matter tensor includes the fluid and electromagnetic contributions;
$n_b$ is the baryon number density. The sign in the last equation means
that emission of an electron neutrino reduces the matter electron fraction.

The angular pressure closure remains Minerbo's maximum-entropy closure
\cite{Minerbo1978}. nuX uses the polynomial approximation
\begin{equation}
 \chi(\xi)=\frac13+
       \frac{\xi^2(6-2\xi+6\xi^2)}{15},\qquad
 \xi=\frac{\sqrt{H_{s\mu}H_s^\mu}}{J_s}.
 \label{eq:minerbo}
\end{equation}
It connects the isotropic and freely streaming limits, $\chi=1/3$ and
$\chi=1$. Selecting WeakLib changes the collision terms, not this angular
closure. In particular, a scalar Fermi--Dirac spectrum for collisions does
not turn M1 into an energy-resolved or angle-resolved transport scheme.

\section{Matter state and spectral approximation}

\subsection{A common equation of state}

EOSX supplies the mass density $\rho$, temperature $T$, electron fraction
$Y_e$, chemical potentials, and nuclear composition used by both the fluid
and the interaction calculation. WeakLibInterp reads and interpolates
opacity tables; it does not replace the hydrodynamic equation of state
\cite{WeakLib,WeakLibInterp}. An EOS table used to generate the opacities
is an offline input, not a second matter model evolved during a simulation.

Temperatures and chemical potentials below are expressed in energy units.
Chemical equilibrium for charged-current reactions requires
\begin{equation}
 \mu_{\nue}^{\eq}=\mu_e+\mu_p-\mu_n,\qquad
 \mu_{\anue}^{\eq}=-\mu_{\nue}^{\eq}.
 \label{eq:chemical_equilibrium}
\end{equation}
The heavy-lepton equilibrium chemical potential is zero. The nucleon
rest-mass difference must be included once in Eq.~\eqref{eq:chemical_equilibrium}.
A common reference subtracted from both nucleon chemical potentials
cancels; subtracting each nucleon's own rest mass does not. This convention
must agree between EOSX and the table generator. An arbitrary shift of the
EOS specific internal energy is not a neutrino chemical-potential correction.

\subsection{Reconstructing a spectrum from grey moments}

For the coupled collision model, the scalar occupation of one neutrino
species is approximated by
\begin{equation}
 f_s(\epsilon)=\frac{1}{\exp(\epsilon/T_{\nu,s}-\eta_{\nu,s})+1},
 \label{eq:fd}
\end{equation}
where $\epsilon$ is comoving neutrino energy. The radiation temperature
$T_{\nu,s}$ and degeneracy parameter $\eta_{\nu,s}$ are determined from
the current radiation moments, not set equal to the matter temperature and
chemical potential. With one helicity state per physical species,
\begin{align}
 n_s&=g_s C\int_0^\infty\epsilon^2f_s(\epsilon)\dd\epsilon,
 & C&=\frac{1}{2\pi^2(\hbar c)^3},\label{eq:number_moment}\\
 J_s&=g_s C\int_0^\infty\epsilon^3f_s(\epsilon)\dd\epsilon.
 \label{eq:energy_moment}
\end{align}
Here $g_{\nue}=g_{\anue}=1$ and $g_{\nux}=4$ count physical species.
Thus $n_s=g_s CT_{\nu,s}^3F_2(\eta_{\nu,s})$ and
$J_s=g_s CT_{\nu,s}^4F_3(\eta_{\nu,s})$, with
$F_k(\eta)=\int_0^\infty x^k/(e^{x-\eta}+1)\dd x$.
For the heavy-lepton group, the reconstructed per-species moments are
$n_{\nux}/4$ and $J_{\nux}/4$.

This assumption gives $0\leq f_s\leq1$, but cannot describe an arbitrary
spectrum. For positive number and energy densities, admissible moments
must satisfy
\begin{equation}
 \frac{n_s^4}{g_s CJ_s^3}\leq\frac{64}{81}
 \label{eq:fd_bound}
\end{equation}
where equality is the zero-temperature filled-sphere limit. Number and
energy floors chosen independently can violate this
condition or imply a mean energy outside the tables. Moreover, bounded
scalar occupations do not establish a complete fermionic angular closure
when the radiation is strongly directed.

An alternative, simpler prescription averages emission, absorption, and
elastic scattering over the matter-equilibrium spectrum. It does not
describe energy redistribution of a general existing radiation field.
The pair, inelastic, and bremsstrahlung coupling discussed below uses the
reconstructed spectrum of Eq.~\eqref{eq:fd} instead.

\section{Included interactions}

The distinction between an interaction kernel and its grey average is
essential. A kernel specifies emission, absorption, or scattering as a
function of neutrino energies. The grey source additionally depends on
the assumed radiation spectrum and, for two-particle processes, on the
partner population \cite{Bruenn1985,Burrows2006,Chiesa2025}.

\subsection{Charged-current emission and absorption}

The electron-flavor reactions on free nucleons are
\begin{align}
 e^-+p&\rightleftharpoons n+\nue,\\
 e^++n&\rightleftharpoons p+\anue.
\end{align}
They exchange radiation energy and electron lepton number with matter.
WeakLib stores a corrected absorption coefficient $\chi_s(\epsilon)$
that includes the stimulated-absorption convention. The occupation source
can therefore be written in local proper time $\tau$ as
\begin{equation}
 \left.\frac{\partial f_s}{\partial\tau}\right|_{\mathrm{cc}}
   =c\chi_s(\epsilon)
       [f_s^{\eq}(\epsilon)-f_s(\epsilon)].
 \label{eq:cc_source}
\end{equation}
Multiplying $\chi_s$ by an additional final-state neutrino blocking factor
would count that correction twice. Matter blocking, thresholds, recoil,
and weak-magnetism corrections depend on the model used to generate the
tables; they are not guaranteed by the name WeakLib.

For a nonempty population, the number- and energy-weighted absorption
coefficients are
\begin{equation}
 \bar\chi_s^{(q)}=
 \frac{\int\epsilon^q\chi_s(\epsilon)f_s(\epsilon)\dd\epsilon}
      {\int\epsilon^q f_s(\epsilon)\dd\epsilon},\qquad q=2,3.
\end{equation}
The corresponding emissivities use $\chi_s f_s^{\eq}$ rather than
$\chi_s f_s$. Emission is therefore finite in an initially empty radiation
field and does not require dividing an emissivity by the evolved density.

The present reference tables use the isoenergetic free-nucleon model.
Additional nuclear electron-capture data are not included in this reference
set, and the separate nuclear-capture table is not coupled. Other WeakLib
tables may already incorporate nuclear capture in their main absorption
arrays. Such a contribution must not be added a second time.

\subsection{Elastic scattering}

Neutral-current scattering,
\begin{equation}
 \nu_s+b\longrightarrow\nu_s+b,\qquad b\in\{n,p\},
\end{equation}
redistributes direction without changing neutrino energy in the
isoenergetic approximation. Depending on the supplied tables, coherent
scattering on nuclei and scattering on alpha particles can also contribute.
The zeroth and first Legendre coefficients are needed for the transport
opacity. With the WeakLib expansion convention used here,
\begin{equation}
 \kappa_{\mathrm{tr},s}(\epsilon)
    =4\pi\epsilon^2
       \left[\phi_{0,s}(\epsilon)-\frac{\phi_{1,s}(\epsilon)}{3}\right].
 \label{eq:transport_opacity}
\end{equation}
The factor $1/3$ follows from the normalization of the stored expansion
coefficients. Elastic scattering damps the comoving flux but contributes
neither a number source nor a comoving energy source. For the combined
heavy-lepton group, neutrino and antineutrino coefficients are averaged;
the opacity itself is not multiplied by four.

\subsection{Electron--positron annihilation and bremsstrahlung}

The pair-production reactions are
\begin{align}
 e^-+e^+&\rightleftharpoons\nu_s+\bar\nu_s,\\
 b+b'&\rightleftharpoons b+b'+\nu_s+\bar\nu_s,
 \qquad b,b'\in\{n,p\}.
\end{align}
The second process is nucleon--nucleon bremsstrahlung. Both create a
neutrino and an antineutrino in one event and conserve electron lepton
number. Bremsstrahlung is especially relevant to heavy-lepton production
in dense matter, where charged-current emission is absent in the reaction
set considered here \cite{HannestadRaffelt1998}.

Writing the angularly averaged production and annihilation kernels as
$\mathcal R^+$ and $\mathcal R^-$, respectively, the pair collision term
contains
\begin{equation}
 \mathcal R^+(1-f_s)(1-f_{\bar s})
       -\mathcal R^- f_s f_{\bar s}.
 \label{eq:pair_source}
\end{equation}
Matter in thermal equilibrium imposes
\begin{equation}
 \mathcal R^+(\epsilon,\epsilon')
    =e^{-(\epsilon+\epsilon')/T}
       \mathcal R^-(\epsilon,\epsilon').
 \label{eq:pair_balance}
\end{equation}
The partner occupations in Eq.~\eqref{eq:pair_source} are reconstructed from
the current antineutrino and neutrino moments. They are not independently
replaced by thermal populations. Consequently, the two species receive
equal net number sources, although their energy sources need not be equal.
The heavy-lepton sum contains two flavor pairs, accounting for four
individual species without duplicating a physical pair event.

The WeakLib coupling uses the Hannestad--Raffelt spin-response model for
bremsstrahlung \cite{HannestadRaffelt1998}. The neutron--neutron,
proton--proton, and mixed contributions are assembled from the supplied
response tables and free-nucleon fractions. This is not the same model as
the Burrows--Reddy--Thompson prescription available in NuRates and the
analytic WeakRates fit \cite{Burrows2006}. Agreement of the reaction names
therefore does not imply agreement of their rates.

\subsection{Inelastic scattering on electrons and positrons}

The reactions
\begin{equation}
 \nu_s+e^\pm\longrightarrow\nu_s+e^\pm
\end{equation}
exchange energy and momentum with matter but preserve neutrino number.
Their kernels connect different neutrino energies. For any interacting
pair of energies, detailed balance requires
\begin{align}
 &\mathcal R(\epsilon'\!\to\epsilon)
       f^{\eq}(\epsilon')[1-f^{\eq}(\epsilon)]\nonumber\\
 &\hspace{1em}=\mathcal R(\epsilon\!\to\epsilon')
       f^{\eq}(\epsilon)[1-f^{\eq}(\epsilon')].
\end{align}
The integrated number source is identically zero, whereas the energy
source can heat or cool the radiation. This channel is important for
heavy-lepton thermalization \cite{Chiesa2025}.

The present pair and inelastic terms retain the angular monopole of the
collision kernels: they use the angle-averaged occupations. Flux attenuation
is evaluated with a spectrum proportional to $f_s(\epsilon)H_s^\mu/J_s$.
This is an additional approximation beyond grey M1, not the complete
angular dependence of the WeakLib kernels. The Minerbo transport closure
is retained independently.

\subsection{Plasmon decay}

Collective electromagnetic excitations of the electron plasma can decay,
\begin{equation}
 \gamma^*\rightleftharpoons\nu_s+\bar\nu_s.
\end{equation}
This contribution is not taken from a native WeakLib table. It is a
separately selectable analytic supplement using the WeakRates plasmon
emission formula and a monoenergetic approximation. Its inverse is
constructed with the same thermal balance relation and the actual partner
blocking factors. Only this channel is borrowed from WeakRates; adding it
does not superpose an additional complete rate calculation.

\section{Grey collision sources and time integration}

Let $\mathcal C_s(\epsilon)$ denote the net angle-averaged occupation
change per unit path length, with the kernel normalizations included.
For the grouped heavy species, $\mathcal C_s$ is the mean of the four
physical species' collision terms. The physical number and energy
production rates are
\begin{align}
 Q_{N,s}&=cg_s C\int\epsilon^2\mathcal C_s(\epsilon)\dd\epsilon,\\
 Q_{E,s}&=cg_s C\int\epsilon^3\mathcal C_s(\epsilon)\dd\epsilon.
\end{align}
Emission and attenuation must be averaged separately with these two
weights; a number opacity is not generally an energy opacity.
The finite-band integrals use a common quadrature in physical energy,
not in its logarithm. The same nodes, weights, and reaction normalizations
are used for forward and inverse processes so that discrete thermal
equilibrium is preserved.

Collision coefficients are recomputed from the trial radiation state
during the coupled nonlinear solve. This includes the partner dependence
of the pair reactions. Matter density, temperature, composition, velocity,
and geometry are held fixed within that local radiation solve; the
matter exchange is applied according to the selected time integrator.
This is not a simultaneous nonlinear solution for the matter equation of
state and all radiation moments.

The available integration choices include an RK4 fluid evolution with a
semi-implicit M1 predictor--corrector and implicit--explicit Runge--Kutta
methods \cite{Izquierdo2023}. RK4 in the former name does not establish
fourth-order accuracy for the combined stiff collision update. For an
implicit--explicit method, the collision increments must enter with its
stage and final weights; an extra unweighted final collision solve would
define a different method. Time-integration errors and spectral-average
errors are separate from errors in the microscopic rates.

\section{Comparison with WeakRates and NuRates}

Table~\ref{tab:reactions} summarizes the reactions available through the
three nuX couplings. It describes capabilities, not a requirement to
activate every channel in every calculation.

\begin{table*}[t]
\caption{Reaction content of the nuX rate couplings. ``Table dependent''
means that the physical process and its corrections must be present in
the supplied WeakLib tables. The WeakLib column describes the coupled
angle-averaged collision model.}
\label{tab:reactions}
\begin{ruledtabular}
\begin{tabular}{llll}
Interaction & WeakRates & NuRates & WeakLib coupling \\
\hline
Free-nucleon charged currents & Analytic grey rates & Energy-dependent rates & Tabulated absorption \\
Elastic scattering on nucleons & Included & Included & Included \\
Elastic scattering on nuclei & Included & Not in the present coupling & Table dependent \\
$e^-e^+$ pair annihilation & Analytic grey emission$^{a}$ & Two-energy kernels & Tabulated two-energy kernels \\
Nucleon bremsstrahlung & Optional analytic fit$^{a}$ & Selectable microscopic models$^{b}$ & Hannestad--Raffelt tables \\
Inelastic $\nu e^\pm$ scattering & Not included & Included & Tabulated kernels \\
Plasmon decay & Analytic emission$^{a}$ & Not included & Optional analytic supplement \\
Nuclear electron capture & Not included & Not included & Not in the present reference set$^{c}$ \\
\end{tabular}
\end{ruledtabular}
\vspace{2pt}
\parbox{\textwidth}{\footnotesize
$^{a}$In the legacy M1 use of WeakRates, electron-flavor emissivities are
reconstructed from charged-current absorption. Electron-flavor pair,
plasmon, and bremsstrahlung contributions require the optional non-beta
detailed-balance extension. Heavy-lepton emission uses these pair sources
directly.
$^{b}$Hannestad--Raffelt and Burrows--Reddy--Thompson prescriptions are
available for GPU calculations. The Guo--Mart\'inez-Pinedo option requires
additional device-table support in the present coupling.
$^{c}$Main absorption arrays may contain nuclear capture if generated with
that physics, but a separate nuclear-capture table is not coupled.}
\end{table*}

WeakRates evaluates analytic expressions with thermal spectral averages
and approximate blocking \cite{WeakRates}. It avoids explicit two-energy
collision integrals, but lacks inelastic neutrino--electron scattering. Its
legacy electron-flavor prescription and the optional extension in
Table~\ref{tab:reactions} must be distinguished when comparing luminosities
or cooling rates.

NuRates uses the \textsc{bns\_nurates} library, which evaluates
energy-dependent rates and two-energy kernels during the calculation
\cite{Chiesa2025,BNSNuRates}. It provides optional mean-field,
effective-mass, weak-magnetism, recoil, and nucleon-decay corrections.
Their availability does not mean they are enabled in a particular run.
The current nuX collision prescription reconstructs a moment-matched
Fermi--Dirac spectrum; this differs from the more general distribution
mixture offered by the underlying library. Inelastic scattering contributes
to energy exchange, not to neutrino-number creation.

WeakLib moves the microscopic rate calculation into an offline table
generation step. It can represent additional nuclear-composition and
scattering corrections when the chosen tables contain them, but still
requires spectral integrals during grey transport. Interpolation error,
energy quadrature, table extent, and the adopted angular and spectral
approximations remain. It is therefore not possible to rank the three
couplings by physical accuracy from their names or number of activated
channels alone. A controlled comparison must match the EOS, composition,
reaction models, corrections, boundary conditions, resolution, numerical
floors, and time-integration method. Adding a reaction can change the
solution physically without demonstrating that either solution is more
accurate.

\section{Using WeakLib in nuX}
\label{sec:usage}

\subsection{What the interpolator supplies}

WeakLibInterp loads the interaction tables and makes their data available
for local interpolation on the GPU \cite{WeakLibInterp}. It evaluates
absorption coefficients, elastic-scattering coefficients, and pair or
inelastic kernels at the requested matter state. nuX then performs the
energy integrals using the inferred radiation spectrum. The interpolator
does not solve radiation transport, choose a neutrino spectrum, or supply
the fluid pressure and internal energy.

Absorption and elastic tables use $\rho$, $T$, and $Y_e$; the present
interpolation requires uniform spacing in $\log\rho$, $\log T$, and $Y_e$.
Pair and neutrino--electron scattering tables instead use the matter
temperature and electron degeneracy $\eta_e=\mu_e/T$ for their thermodynamic
coordinates. The supported pair and inelastic tables require positive
$\eta_e$ within their tabulated range. The temperature passed to the tables
is converted from MeV to kelvin. Stored logarithms, additive offsets, units,
and angular normalizations must be interpreted as specified by the table
family; a signed first Legendre coefficient is not an invalid negative
physical opacity.

The present nuX quadrature queries the coefficients at stored energy
nodes. All enabled channels must share the same positive, increasing
physical energy grid. The nodes need not be uniformly spaced in energy.
Coefficients are interpolated in their matter coordinates, whereas the
subsequent quadrature uses physical energy separations. No out-of-range
opacity extrapolation is used by this coupling.

\subsection{Preparing compatible tables}

Existing tables may be used if they describe the same matter model and
cover the required state space. Otherwise they must be generated with
WeakLib's opacity-table generator \cite{WeakLibTables}. The preparation
has four steps:
\begin{enumerate}
 \item Supply an EOS in the generator's input format. If the original
 EOSX table uses another format, convert it offline, preserving density,
 temperature, chemical potentials, and the free-nucleon and nuclear mass
 fractions. Check the decoded quantities against the original EOS and
 verify Eq.~\eqref{eq:chemical_equilibrium}. A change of file format must
 not change the physical EOS or silently move unrepresented nuclear
 species into other composition fractions.

 \item Choose the reaction model and grids in the table-creation driver.
 Specify the charged-current approximation, scattering corrections,
 energy nodes, thermodynamic extent, and electron-degeneracy range.
 Generate the emission/absorption, elastic, pair, inelastic, and
 bremsstrahlung families with compatible energy grids. This is where
 microscopic corrections are selected; interpolation cannot add a
 correction absent from a table.

 \item Configure a Fortran compiler and HDF5 for the WeakLib generator,
 build its \texttt{wlCreateOpacityTable} target, and run the resulting
 executable with the prepared EOS selected by the driver. Use a fresh
 output directory. This produces the channel-specific HDF5 tables;
 no native plasmon table is required for the analytic supplement.
 The generated EOS-format copy is not used as the runtime matter EOS.

 \item Check table units and grid coverage, compare interpolated rates
 with direct generator values, and test reaction balance before an
 evolution. Repeat representative collision calculations with finer
 energy and thermodynamic grids and a wider energy interval. Small
 omitted number and energy tails alone do not bound errors in an
 opacity-weighted integral.
\end{enumerate}

Bremsstrahlung needs an additional range check: the response is sampled
at effective densities $\rho X_p$, $\rho X_n$, and
$\rho\sqrt{X_pX_n}$. These can be much smaller than the total mass density.
The bremsstrahlung table must cover the positive effective densities that
occur in the calculation, not just the total-density range of the EOS.
Similarly, dilute radiation and strongly degenerate radiation can require
energy coverage different from an equilibrium spectrum at the matter
temperature. A table generated for one EOS, such as SFHo, is not a
consistent substitute for a missing APRLDP table.

\subsection{Selecting the coupled reaction model}

The executable must include WeakLibInterp and the nuX WeakLib support.
Activate \texttt{nuX\_WeakLib} together with EOSX and M1 in the simulation;
the interpolator is a supporting library, not a second EOS to activate.
The following fragment selects the coupled collision model and its
available channels. Replace the example paths with compatible tables.

\begingroup\footnotesize
\begin{verbatim}
nuX_M1::rates_lib = "WeakLib"
nuX_M1::opacity_tau_trap = -1
nuX_M1::opacity_corr_fac_max = 1
nuX_M1::source_limiter = -1

nuX_WeakLib::collision_model = "monopole"
nuX_WeakLib::include_emab = yes
nuX_WeakLib::include_iso = yes
nuX_WeakLib::include_pair = yes
nuX_WeakLib::include_nes = yes
nuX_WeakLib::include_brem = yes
nuX_WeakLib::include_plasmon = yes
nuX_WeakLib::include_nuclei_ec = no

nuX_WeakLib::emab_table = "tables/EmAb.h5"
nuX_WeakLib::iso_table = "tables/Iso.h5"
nuX_WeakLib::pair_table = "tables/Pair.h5"
nuX_WeakLib::nes_table = "tables/NES.h5"
nuX_WeakLib::brem_table = "tables/Brem.h5"

nuX_WeakLib::chemical_potentials = "total"
nuX_WeakLib::spectral_tail_tolerance = 1e-2
nuX_WeakLib::run_self_tests = yes
\end{verbatim}
\endgroup

The disabled trapped-equilibrium blend and unit mean-energy correction
prevent an additional equilibrium prescription from being superposed on
the explicitly reconstructed collision spectrum. Disabling the source
limiter retains the accepted nonlinear source solution. Radiation floors
still require a physically consistent choice of number and energy density.

The scattering table also requires a convention declaration. For files
that explicitly label the raw kernels in
$\mathrm{cm}^{-1}\,\mathrm{MeV}^{-2}$, retain the default
\texttt{iso\_table\_convention = "explicit-kernel-units"}.
For native FullWeakPhysics tables verified to use the corresponding raw
kernel normalization despite their older unit label, select
\texttt{"weaklib-fullweakphysics-5836de98"} instead. This declaration
does not alter the rates or establish EOS compatibility; it identifies
how to interpret the supplied scattering coefficients.

The chemical-potential choice shown above assumes the total-potential
convention of Eq.~\eqref{eq:chemical_equilibrium}. The alternative
\texttt{"rest-mass-subtracted"} choice is appropriate only if the EOS
removes each nucleon's own rest mass. The spectral-tail
tolerance is a coverage check, not an accuracy guarantee for the rates.
For validation, first enable the collision self-tests and run a small
thermal-equilibrium and relaxation problem. Only then compare full
evolutions with the same matter state, enabled reactions, and transport
settings. These checks are necessary but do not replace convergence
studies in space, time, and the tabulated energy grid.

\begin{thebibliography}{99}

\bibitem{Minerbo1978}
G.~N.~Minerbo, Maximum entropy Eddington factors,
\href{https://doi.org/10.1016/0022-4073(78)90024-9}{J. Quant. Spectrosc.
Radiat. Transfer \textbf{20}, 541 (1978)}.

\bibitem{Bruenn1985}
S.~W.~Bruenn, Stellar core collapse: Numerical model and infall epoch,
\href{https://doi.org/10.1086/191056}{Astrophys. J. Suppl. Ser.
\textbf{58}, 771 (1985)}.

\bibitem{HannestadRaffelt1998}
S.~Hannestad and G.~Raffelt, Supernova neutrino opacity from
nucleon--nucleon bremsstrahlung and related processes,
\href{https://doi.org/10.1086/306303}{Astrophys. J. \textbf{507}, 339 (1998)},
\href{https://arxiv.org/abs/astro-ph/9711132}{arXiv:astro-ph/9711132}.

\bibitem{Burrows2006}
A.~Burrows, S.~Reddy, and T.~A.~Thompson, Neutrino opacities in nuclear matter,
\href{https://doi.org/10.1016/j.nuclphysa.2004.06.012}{Nucl. Phys. A
\textbf{777}, 356 (2006)},
\href{https://arxiv.org/abs/astro-ph/0404432}{arXiv:astro-ph/0404432}.

\bibitem{Chiesa2025}
L.~Chiesa, M.~Bhattacharyya, F.~Mazzini, F.~M.~Guercilena, A.~Perego,
and D.~Radice, Open-source library for performance-portable neutrino
reaction rates: Application to neutron star mergers,
\href{https://doi.org/10.1103/PhysRevD.111.063053}{Phys. Rev. D
\textbf{111}, 063053 (2025)},
\href{https://arxiv.org/abs/2412.04570}{arXiv:2412.04570}.

\bibitem{MusolinoRezzolla2023}
C.~Musolino and L.~Rezzolla, A practical guide to a moment approach for
neutrino transport in numerical relativity,
\href{https://arxiv.org/abs/2304.09168}{arXiv:2304.09168}.

\bibitem{Izquierdo2023}
M.~R.~Izquierdo, L.~Pareschi, B.~Mi\~nano, J.~Mass\'o, and C.~Palenzuela,
Global high-order numerical schemes for the time evolution of the
general relativistic radiation magnetohydrodynamics equations,
\href{https://doi.org/10.1088/1361-6382/acd97f}{Class. Quantum Grav. (2023)},
\href{https://arxiv.org/abs/2211.00027}{arXiv:2211.00027}.

\bibitem{WeakLib}
WeakLib, neutrino opacity and equation-of-state tables,
\url{https://github.com/starkiller-astro/weaklib}.

\bibitem{WeakLibInterp}
WeakLibInterp, interpolation of WeakLib interaction tables,
\url{https://github.com/jaykalinani/WeakLibInterp}.

\bibitem{BNSNuRates}
\textsc{bns\_nurates}, neutrino interaction-rate library,
\url{https://github.com/RelNucAs/bns_nurates}.

\bibitem{WeakRates}
WeakRates analytic neutrino interaction rates as used in nuX,
\url{https://github.com/jaykalinani/nuX/tree/odesolvers/nuX_WeakRates}.

\bibitem{WeakLibTables}
WeakLib, opacity-table generation guide,
\url{https://github.com/starkiller-astro/weaklib/blob/master/TableCreator/OpacitySource/BuildOpacityTable.md}.

\end{thebibliography}

\end{document}
